\documentclass[11pt]{article}

\usepackage[margin=0.8in]{geometry}
\usepackage{amsmath,amssymb}
\usepackage{enumitem}
\usepackage{xcolor}
\usepackage{tcolorbox}
\usepackage{fancyhdr}
\usepackage{lastpage}
\usepackage{hyperref}
\usepackage{microtype}

\definecolor{uwpurple}{HTML}{4B2E83}
\definecolor{uwgold}{HTML}{B7A57A}
\definecolor{softpurple}{HTML}{F4F1FA}
\definecolor{softgold}{HTML}{FBF7EA}
\definecolor{deepgold}{HTML}{8A6F2A}

\tcbuselibrary{breakable,skins}

\pagestyle{fancy}
\fancyhf{}
\lhead{\textbf{MATH 394: Probability I}}
\rhead{\textbf{Makeup Midterm Solutions}}
\cfoot{\thepage\ of \pageref{LastPage}}
\setlength{\headheight}{14pt}

\setlength{\parindent}{0pt}
\setlength{\parskip}{0.45em}

\newcommand{\course}{MATH 394: Probability I}
\newcommand{\instructor}{Arman Jahangiri}
\newcommand{\department}{Department of Mathematics}
\newcommand{\university}{University of Washington}

\newtcolorbox{problemheader}{
    colback=softpurple,
    colframe=uwpurple,
    boxrule=0.8pt,
    arc=2mm,
    left=2mm,
    right=2mm,
    top=1mm,
    bottom=1mm
}

\newtcolorbox{solutionbox}[1][] {
    breakable,
    enhanced,
    colback=softgold,
    colframe=deepgold,
    coltitle=white,
    colbacktitle=uwpurple,
    fonttitle=\bfseries,
    title={Solution #1},
    boxrule=0.85pt,
    arc=2mm,
    left=2.5mm,
    right=2.5mm,
    top=1.5mm,
    bottom=1.5mm,
    attach boxed title to top left={xshift=2.5mm,yshift=-2mm},
    boxed title style={
        colback=uwpurple,
        colframe=uwpurple,
        arc=1.5mm,
        boxrule=0pt,
        left=1.6mm,
        right=1.6mm,
        top=0.6mm,
        bottom=0.6mm
    }
}

\newtcolorbox{finalanswerbox}[1][] {
    colback=white,
    colframe=uwpurple,
    boxrule=0.75pt,
    arc=2mm,
    left=2mm,
    right=2mm,
    top=1mm,
    bottom=1mm,
    title={\textbf{Final Answer #1}},
    coltitle=uwpurple,
    colbacktitle=softpurple,
    fonttitle=\bfseries
}

\begin{document}

% ============================================================
% TITLE PAGE
% ============================================================

\begin{titlepage}
\centering
\vspace*{1.2cm}

{\Large \textbf{\university}}\\
{\large \department}

\vspace{1.4cm}

{\Huge \color{uwpurple}\textbf{Makeup Midterm  Solutions}}\\[0.25cm]
{\Large \textbf{\course}}\\[0.25cm]
{\large Summer 2026 \quad | \quad Instructor: \instructor}

\vspace{1cm}

\begin{tcolorbox}[
    width=0.78\textwidth,
    colback=softpurple,
    colframe=uwpurple,
    boxrule=0.9pt,
    arc=3mm
]
\centering

\textbf{Answering time:} 60 minutes

\textbf{Total:} 60 points

\end{tcolorbox}

\end{titlepage}


% ============================================================
% PROBLEM 1
% ============================================================

\begin{problemheader}
\textbf{Problem 1.} \hfill \textbf{10 points}
\end{problemheader}

A box contains 6 red balls, 5 blue balls, and 4 green balls. Five balls
are selected uniformly at random without replacement.

Find the probability that the selection contains exactly 2 red balls,
exactly 2 blue balls, and exactly 1 green ball.

Give an exact answer. A correct expression involving binomial coefficients
is acceptable.

\begin{solutionbox}

There are
\[
\binom{15}{5}
\]
possible selections of 5 balls from the 15 balls, and all such selections
are equally likely.

To obtain the required composition, we choose

\begin{itemize}[leftmargin=*]
    \item 2 of the 6 red balls,
    \item 2 of the 5 blue balls,
    \item 1 of the 4 green balls.
\end{itemize}

Therefore, the number of favorable selections is
\[
\binom{6}{2}\binom{5}{2}\binom{4}{1}.
\]

Hence,
\[
P(\text{required composition})
=
\frac{
\binom{6}{2}\binom{5}{2}\binom{4}{1}
}{
\binom{15}{5}
}.
\]

Evaluating,
\[
\binom{6}{2}\binom{5}{2}\binom{4}{1}
=
15\cdot 10\cdot 4
=
600,
\]
while
\[
\binom{15}{5}=3003.
\]

Thus,
\[
P(\text{required composition})
=
\frac{600}{3003}
=
\frac{200}{1001}.
\]

\begin{finalanswerbox}
\[
P(\text{exactly 2 red, 2 blue, and 1 green})
=
\frac{\binom62\binom52\binom41}{\binom{15}{5}}
=
\frac{200}{1001}.
\]
\end{finalanswerbox}

\end{solutionbox}


\newpage

% ============================================================
% PROBLEM 2
% ============================================================

\begin{problemheader}
\textbf{Problem 2.} \hfill \textbf{10 points}
\end{problemheader}

A company purchases computer chips from three suppliers.

Supplier $A$ provides $50\%$ of the chips, Supplier $B$ provides $30\%$,
and Supplier $C$ provides the remaining $20\%$.

The probabilities that a chip is defective are
\[
P(D\mid A)=0.01,\qquad
P(D\mid B)=0.03,\qquad
P(D\mid C)=0.06,
\]
where $D$ denotes the event that a chip is defective.

A chip is selected at random and is found to be defective.

\begin{enumerate}[label=(\alph*),leftmargin=*,itemsep=0.8em]

    \item Find the probability that the defective chip came from
    Supplier $B$.
    \hfill \textbf{(6 points)}

    \item Find the probability that the defective chip came from
    either Supplier $B$ or Supplier $C$.
    \hfill \textbf{(4 points)}

\end{enumerate}


\begin{solutionbox}

We are given
\[
P(A)=0.50,\qquad
P(B)=0.30,\qquad
P(C)=0.20,
\]
and
\[
P(D\mid A)=0.01,\qquad
P(D\mid B)=0.03,\qquad
P(D\mid C)=0.06.
\]

Since $A$, $B$, and $C$ form a partition according to the supplier,
the law of total probability gives
\begin{align*}
P(D)
&=
P(D\mid A)P(A)
+
P(D\mid B)P(B)
+
P(D\mid C)P(C)\\
&=
(0.01)(0.50)
+
(0.03)(0.30)
+
(0.06)(0.20)\\
&=
0.005+0.009+0.012\\
&=
0.026.
\end{align*}


\textbf{(a)}

By Bayes' rule,
\begin{align*}
P(B\mid D)
&=
\frac{P(D\mid B)P(B)}{P(D)}\\
&=
\frac{(0.03)(0.30)}{0.026}\\
&=
\frac{0.009}{0.026}\\
&=
\frac{9}{26}.
\end{align*}

Thus,
\[
P(B\mid D)
=
\frac{9}{26}
\approx 0.3462.
\]


\textbf{(b)}

Since the suppliers are mutually exclusive,
\[
P(B\cup C\mid D)
=
P(B\mid D)+P(C\mid D).
\]

We may compute this directly:
\begin{align*}
P(B\cup C\mid D)
&=
\frac{
P(D\mid B)P(B)+P(D\mid C)P(C)
}{
P(D)
}\\
&=
\frac{
(0.03)(0.30)+(0.06)(0.20)
}{
0.026
}\\
&=
\frac{0.009+0.012}{0.026}\\
&=
\frac{0.021}{0.026}\\
&=
\frac{21}{26}.
\end{align*}

Therefore,
\[
P(B\cup C\mid D)
=
\frac{21}{26}
\approx0.8077.
\]

\begin{finalanswerbox}
\[
\text{(a)}\qquad
P(B\mid D)=\frac{9}{26}\approx0.3462,
\]
\[
\text{(b)}\qquad
P(B\cup C\mid D)=\frac{21}{26}\approx0.8077.
\]
\end{finalanswerbox}

\end{solutionbox}


\newpage

% ============================================================
% PROBLEM 3
% ============================================================

\begin{problemheader}
\textbf{Problem 3.} \hfill \textbf{10 points}
\end{problemheader}

Let $A$ and $B$ be events.

\begin{enumerate}[label=(\alph*),leftmargin=*,itemsep=1em]

    \item Suppose
    \[
    P(A)=0.7,
    \qquad
    P(B)=0.6.
    \]
    What is the smallest possible value of $P(A\cap B)$?
    Justify your answer.
    \hfill \textbf{(3 points)}

    \item Suppose $A$ and $B$ are disjoint and
    \[
    P(A)>0,
    \qquad
    P(B)>0.
    \]
    Can $A$ and $B$ be independent? Justify your answer using the
    definition of independence.
    \hfill \textbf{(3 points)}

    \item Suppose $A$ and $B$ are independent. Prove that
    $A$ and $B^c$ are also independent.
    \hfill \textbf{(4 points)}

\end{enumerate}


\begin{solutionbox}

\textbf{(a)}

Using
\[
P(A\cup B)
=
P(A)+P(B)-P(A\cap B),
\]
and the fact that
\[
P(A\cup B)\leq 1,
\]
we obtain
\[
P(A)+P(B)-P(A\cap B)\leq1.
\]

Therefore,
\[
P(A\cap B)
\geq
P(A)+P(B)-1.
\]

Substituting the given values,
\[
P(A\cap B)
\geq
0.7+0.6-1
=
0.3.
\]

Thus the smallest possible value is
\[
P(A\cap B)=0.3.
\]

This value is possible, for example, when
\[
P(A\cup B)=1.
\]


\textbf{(b)}

If $A$ and $B$ are disjoint, then
\[
A\cap B=\emptyset,
\]
so
\[
P(A\cap B)=0.
\]

However, if $A$ and $B$ were independent, the definition of independence
would require
\[
P(A\cap B)
=
P(A)P(B).
\]

Since
\[
P(A)>0
\qquad\text{and}\qquad
P(B)>0,
\]
we have
\[
P(A)P(B)>0.
\]

Thus independence would imply
\[
P(A\cap B)>0,
\]
which contradicts
\[
P(A\cap B)=0.
\]

Therefore, two disjoint events having positive probability cannot be
independent.


\textbf{(c)}

We want to prove that
\[
P(A\cap B^c)
=
P(A)P(B^c).
\]

Notice that $A$ is the disjoint union
\[
A
=
(A\cap B)\cup(A\cap B^c).
\]

Therefore,
\[
P(A)
=
P(A\cap B)+P(A\cap B^c),
\]
and hence
\[
P(A\cap B^c)
=
P(A)-P(A\cap B).
\]

Since $A$ and $B$ are independent,
\[
P(A\cap B)=P(A)P(B).
\]

Thus,
\begin{align*}
P(A\cap B^c)
&=
P(A)-P(A)P(B)\\
&=
P(A)\bigl(1-P(B)\bigr)\\
&=
P(A)P(B^c).
\end{align*}

Therefore, by the definition of independence,
\[
A\quad\text{and}\quad B^c
\]
are independent.

\begin{finalanswerbox}

\textbf{(a)}
\[
\min P(A\cap B)=0.3.
\]

\textbf{(b)} No. Disjoint events with positive probabilities cannot be
independent.

\textbf{(c)}
\[
P(A\cap B^c)=P(A)P(B^c),
\]
so $A$ and $B^c$ are independent.

\end{finalanswerbox}

\end{solutionbox}


\newpage

% ============================================================
% PROBLEM 4
% ============================================================

\begin{problemheader}
\textbf{Problem 4.} \hfill \textbf{15 points}
\end{problemheader}

Suppose that $X$ has probability density function
\[
f_X(x)=
\begin{cases}
c(1+x), & 0<x<2,\\
0, & \text{otherwise},
\end{cases}
\]
where $c$ is a constant.

\begin{enumerate}[label=(\alph*),leftmargin=*,itemsep=0.8em]

    \item Find $c$.
    \hfill \textbf{(3 points)}

    \item Find the cumulative distribution function $F_X(x)$
    for all real $x$.
    \hfill \textbf{(5 points)}

    \item Find
    \[
    P\left(\frac12<X<\frac32\right).
    \]
    \hfill \textbf{(2 points)}

    \item Find $\mathbb{E}[X]$ and $\operatorname{Var}(X)$.
    \hfill \textbf{(5 points)}

\end{enumerate}


\begin{solutionbox}

\textbf{(a)}

A probability density function must integrate to $1$. Therefore,
\begin{align*}
1
&=
\int_{-\infty}^{\infty} f_X(x)\,dx\\
&=
c\int_0^2(1+x)\,dx\\
&=
c\left[x+\frac{x^2}{2}\right]_0^2\\
&=
c(2+2)\\
&=
4c.
\end{align*}

Hence,
\[
c=\frac14.
\]


\textbf{(b)}

For $x\leq0$,
\[
F_X(x)=0.
\]

For $0<x<2$,
\begin{align*}
F_X(x)
&=
\int_{-\infty}^{x}f_X(t)\,dt\\
&=
\int_0^x\frac14(1+t)\,dt\\
&=
\frac14
\left[t+\frac{t^2}{2}\right]_0^x\\
&=
\frac{x}{4}+\frac{x^2}{8}\\
&=
\frac{x^2+2x}{8}.
\end{align*}

For $x\geq2$,
\[
F_X(x)=1.
\]

Therefore,
\[
F_X(x)
=
\begin{cases}
0, & x\leq0,\\[2mm]
\displaystyle \frac{x^2+2x}{8}, & 0<x<2,\\[3mm]
1, & x\geq2.
\end{cases}
\]


\textbf{(c)}

Using the CDF,
\begin{align*}
P\left(\frac12<X<\frac32\right)
&=
F_X\left(\frac32\right)
-
F_X\left(\frac12\right).
\end{align*}

Now,
\begin{align*}
F_X\left(\frac32\right)
&=
\frac{
(3/2)^2+2(3/2)
}{8}\\
&=
\frac{\frac94+3}{8}\\
&=
\frac{21}{32},
\end{align*}
while
\begin{align*}
F_X\left(\frac12\right)
&=
\frac{
(1/2)^2+2(1/2)
}{8}\\
&=
\frac{\frac14+1}{8}\\
&=
\frac{5}{32}.
\end{align*}

Therefore,
\[
P\left(\frac12<X<\frac32\right)
=
\frac{21}{32}-\frac{5}{32}
=
\frac{16}{32}
=
\frac12.
\]


\textbf{(d)}

First,
\begin{align*}
\mathbb{E}[X]
&=
\int_{-\infty}^{\infty}x f_X(x)\,dx\\
&=
\frac14\int_0^2x(1+x)\,dx\\
&=
\frac14\int_0^2(x+x^2)\,dx\\
&=
\frac14
\left[
\frac{x^2}{2}+\frac{x^3}{3}
\right]_0^2\\
&=
\frac14
\left(
2+\frac83
\right)\\
&=
\frac14\cdot\frac{14}{3}\\
&=
\frac76.
\end{align*}

Next,
\begin{align*}
\mathbb{E}[X^2]
&=
\int_{-\infty}^{\infty}x^2f_X(x)\,dx\\
&=
\frac14\int_0^2x^2(1+x)\,dx\\
&=
\frac14\int_0^2(x^2+x^3)\,dx\\
&=
\frac14
\left[
\frac{x^3}{3}+\frac{x^4}{4}
\right]_0^2\\
&=
\frac14
\left(
\frac83+4
\right)\\
&=
\frac14\cdot\frac{20}{3}\\
&=
\frac53.
\end{align*}

Therefore,
\begin{align*}
\operatorname{Var}(X)
&=
\mathbb{E}[X^2]
-
\bigl(\mathbb{E}[X]\bigr)^2\\
&=
\frac53
-
\left(\frac76\right)^2\\
&=
\frac{60}{36}
-
\frac{49}{36}\\
&=
\frac{11}{36}.
\end{align*}


\begin{finalanswerbox}

\[
c=\frac14.
\]

\[
F_X(x)
=
\begin{cases}
0, & x\leq0,\\[1mm]
\dfrac{x^2+2x}{8}, & 0<x<2,\\[2mm]
1, & x\geq2.
\end{cases}
\]

\[
P\left(\frac12<X<\frac32\right)=\frac12,
\]

\[
\mathbb{E}[X]=\frac76,
\qquad
\operatorname{Var}(X)=\frac{11}{36}.
\]

\end{finalanswerbox}

\end{solutionbox}

 



\newpage

% ============================================================
% PROBLEM 5
% ============================================================

\begin{problemheader}
\textbf{Problem 5.} \hfill \textbf{15 points}
\end{problemheader}

A factory produces two types of electronic components. Each component is
independently classified as Type I with probability \(0.5\) and Type II
with probability \(0.5\).

Suppose that \(60\) components are produced.

Let \(X\) denote the number of Type I components among the \(60\)
components.

\begin{enumerate}[label=(\alph*),leftmargin=*,itemsep=0.9em]

    \item State the distribution of \(X\).

    Determine whether a normal approximation is appropriate by verifying
    the required conditions. Then state the approximating normal
    distribution, including its mean and variance.

    \hfill \textbf{(6 points)}

    \item Using the normal approximation with an appropriate continuity
    correction, approximate the probability that among the \(60\)
    components there are \textbf{at least \(27\) components of each type}.

    \hfill \textbf{(9 points)}

\end{enumerate}


\begin{solutionbox}

\textbf{(a)}

Each of the \(60\) components is independently Type I with probability
\[
p=0.5.
\]

Therefore,
\[
X\sim\operatorname{Bin}(60,0.5).
\]

To determine whether the normal approximation is appropriate, we check
\[
np=60(0.5)=30
\]
and
\[
n(1-p)=60(0.5)=30.
\]

Since both quantities are greater than \(10\), the normal approximation is
appropriate.

The mean is
\[
\mu=np=30,
\]
and the variance is
\[
\sigma^2=np(1-p)
=
60(0.5)(0.5)
=
15.
\]

Thus, we approximate \(X\) by
\[
Y\sim N(30,15),
\]
where
\[
\sigma=\sqrt{15}.
\]


\textbf{(b)}

If \(X\) is the number of Type I components, then the number of Type II
components is
\[
60-X.
\]

Having at least \(27\) components of each type means
\[
X\geq27
\]
and
\[
60-X\geq27.
\]

The second inequality is equivalent to
\[
X\leq33.
\]

Therefore, the desired event is
\[
27\leq X\leq33.
\]

Using the continuity correction,
\[
P(27\leq X\leq33)
\approx
P(26.5<Y<33.5).
\]

Standardizing,
\begin{align*}
P(26.5<Y<33.5)
&=
P\left(
\frac{26.5-30}{\sqrt{15}}
<
Z
<
\frac{33.5-30}{\sqrt{15}}
\right)\\
&=
P\left(
-\frac{3.5}{\sqrt{15}}
<
Z
<
\frac{3.5}{\sqrt{15}}
\right).
\end{align*}

Since
\[
\frac{3.5}{\sqrt{15}}
\approx0.90,
\]
we obtain
\[
P(27\leq X\leq33)
\approx
P(-0.90<Z<0.90).
\]

Using the standard normal table,
\[
\Phi(0.90)=0.8159,
\]
and
\[
\Phi(-0.90)=0.1841.
\]

Therefore,
\[
P(27\leq X\leq33)
\approx
0.8159-0.1841
=
0.6318.
\]


\begin{finalanswerbox}

\[
X\sim\operatorname{Bin}(60,0.5).
\]

Since
\[
np=n(1-p)=30>10,
\]
the normal approximation is appropriate:
\[
X\approx Y,
\qquad
Y\sim N(30,15).
\]

Finally,
\[
P(\text{at least 27 of each type})
=
P(27\leq X\leq33)
\approx
0.632.
\]

\end{finalanswerbox}

\end{solutionbox}




\end{document}